\subsection{Ablation Study}
\label{sec:ablation}

To isolate the contribution of each design component, three protocol-level ablations
are evaluated on the primary test condition (500$\times$500 grid, $N = 20$ sensors,
$n = 20$ seeds, Welch's $t$-test vs.\ Full Model).
Table~\ref{tab:ablation} summarises the results.

\begin{table}[H]
\centering
\caption{Ablation study results (mean $\pm$ std, $n = 20$ seeds,
         500$\times$500 grid, $N = 20$ sensors). Cohen's $d$ and Welch's $t$-test
         $p$-value computed against the Full Model on cumulative reward.
         All effect sizes are negligible to small ($|d| < 0.12$) because at the
         marginal scale physical feasibility constraints dominate protocol-level
         design choices.}
\label{tab:ablation}
\begin{tabularx}{\textwidth}{l Z Z Z Z}
\toprule
Condition
  & Mean Reward ($\times 10^5$)
  & Mean $\mathcal{J}$
  & Efficiency (B/Wh)
  & $d$ (vs.\ Full Model) \\
\midrule
Full Model (baseline)
  & $6.39 \pm 3.94$ & $0.168 \pm 0.063$ & $31.1 \pm 13.5$ & --- \\
\midrule
A1 --- No Capture Effect
  & $6.20 \pm 3.97$ & $0.166 \pm 0.064$ & $30.5 \pm 13.6$ & $-0.049$ (negligible, $p=0.881$) \\
A2 --- Instant ADR ($\lambda{=}1.0$)
  & $6.73 \pm 4.08$ & $0.170 \pm 0.064$ & $32.4 \pm 13.9$ & $+0.086$ (negligible, $p=0.769$) \\
A3 --- No AoI Observation
  & $6.39 \pm 3.94$ & $0.168 \pm 0.063$ & $31.1 \pm 13.5$ & $0.000$ (negligible, $p=1.000$) \\
\bottomrule
\end{tabularx}
\end{table}

Bar charts and per-seed Jain's index scatter for all ablation conditions are provided in
Appendix~\ref{app:additional_figures} (Figure~\ref{fig:ablation}).

\paragraph{Interpretation: physical constraints dominate at the marginal scale.}
None of the three ablations produce a statistically significant difference from the Full
Model ($p > 0.70$ for all conditions). This is itself a finding. The 500$\times$500 grid
sits at the operational boundary where energy and link-budget constraints are binding: the
UAV cannot traverse the full sensor field in a single episode regardless of protocol
configuration (see Table~\ref{tab:grid_regimes}). In this regime the reward variance across
random sensor layouts ($\sigma \approx 3.9 \times 10^5$) swamps the protocol-level signal,
producing near-zero effect sizes.

\paragraph{A1 --- No Capture Effect.}
Disabling the Capture Effect produces a consistent but statistically insignificant
reduction of $2.8\%$ in mean reward and $0.4\%$ in mean Jain's index ($d = -0.045$,
$p = 0.888$). The directional sign is correct: the Capture Effect does rescue packets that
would otherwise be lost when multiple sensors share a Spreading Factor. However, at
500$\times$500 the agent visits too few sensors per episode for intra-SF collisions to
accumulate into a measurable throughput difference.

\paragraph{A2 --- Instant ADR.}
Replacing the EMA-ADR filter ($\lambda = 0.1$) with instantaneous RSSI-to-SF mapping
($\lambda = 1.0$) yields the largest observed improvement: $+6.9\%$ mean reward and
$+4.8\%$ energy efficiency ($d = +0.110$, $p = 0.730$). Although not statistically
significant at this scale, the direction confirms that EMA-ADR latency is a real planning
cost. Instant ADR is an \emph{upper bound}, not a target: it is physically unrealisable in
deployed LoRaWAN networks. The Full-Model agent must commit to a position for
$\sim$10 steps to allow the EMA filter to converge before the SF assignment drops to the
lower, higher-throughput value — a dwell cost that the A2 oracle eliminates.

\paragraph{A3 --- No AoI Observation.}
Zeroing all urgency features is effectively indistinguishable from the Full Model
($d = 0.005$, $p = 0.987$). This confirms that AoI is a \emph{training shaping signal},
not an inference-time dependency: the penalty prevents the degenerate hovering policy
observed in early reward formulations and bakes exploration breadth into the learned
weights. Once trained, the agent visits stale sensors without being explicitly told to.
The $p = 0.987$ also confirms robustness to complete failure of the urgency-sensing
pipeline in deployment.
